\documentclass[12pt,a4paper]{article}

% ---------------- PACOTES ----------------
\usepackage[utf8]{inputenc}
\usepackage[brazil]{babel}
\usepackage{setspace}
\usepackage{geometry}
\usepackage{indentfirst}
\usepackage{hyperref}

\geometry{a4paper, margin=2.5cm}
\onehalfspacing
\setlength{\parindent}{1.25cm}

% ---------------- TÍTULO ----------------
\title{\textbf{SUS Mais: Sistema de Agendamento de Consultas Online}}
\author{
Grupo de Desenvolvimento -- Código de Alta Performance \\
Instituição de Ensino (substituir) \\
}
\date{\today}

\begin{document}

\maketitle
\thispagestyle{empty}

\newpage

\tableofcontents
\newpage

% ---------------- INTRODUÇÃO ----------------
\section{Introdução}

O acesso a consultas médicas no Sistema Único de Saúde (SUS) ainda apresenta desafios relacionados à alta demanda, filas presenciais e demora no processo de agendamento.

Diante desse cenário, o sistema \textbf{SUS Mais} foi desenvolvido como uma plataforma web voltada para o agendamento online de consultas médicas, com o objetivo de facilitar o acesso da população aos serviços de saúde e reduzir a necessidade de deslocamento físico até unidades de atendimento.

% ---------------- PROBLEMA ----------------
\section{Problema}

O modelo tradicional de agendamento de consultas no SUS apresenta diversas limitações:

\begin{itemize}
    \item Necessidade de comparecimento presencial para marcação;
    \item Longas filas de espera;
    \item Sobrecarga de atendentes nas unidades de saúde;
    \item Dificuldade de acesso para pessoas com mobilidade reduzida;
    \item Ausência de um sistema digital centralizado e eficiente.
\end{itemize}

Esses fatores impactam negativamente a eficiência do atendimento público de saúde.

% ---------------- OBJETIVO ----------------
\section{Objetivo}

O objetivo do SUS Mais é desenvolver um sistema web capaz de:

\begin{itemize}
    \item Permitir agendamento online de consultas médicas;
    \item Disponibilizar especialidades médicas de forma organizada;
    \item Exibir horários disponíveis em tempo real;
    \item Reduzir filas presenciais nas unidades de saúde;
    \item Melhorar a experiência do usuário no acesso à saúde pública.
\end{itemize}

% ---------------- DESCRIÇÃO DO SISTEMA ----------------
\section{Descrição do Sistema}

O SUS Mais é uma aplicação web desenvolvida para simular e facilitar o processo de agendamento de consultas médicas.

O sistema permite que o usuário:

\begin{itemize}
    \item Crie uma conta ou realize login;
    \item Acesse a plataforma de agendamento;
    \item Escolha a unidade de saúde desejada;
    \item Selecione a especialidade médica (como cardiologia, dermatologia e clínica geral);
    \item Visualize horários disponíveis;
    \item Confirme o agendamento da consulta.
\end{itemize}

A proposta é oferecer uma interface simples, intuitiva e acessível para todos os usuários.

% ---------------- FUNCIONAMENTO ----------------
\section{Funcionamento do Sistema}

O fluxo de funcionamento do sistema SUS Mais ocorre da seguinte forma:

\begin{enumerate}
    \item O usuário acessa a aplicação web;
    \item Realiza autenticação (login/cadastro);
    \item Seleciona a unidade de saúde;
    \item Escolhe a especialidade médica;
    \item Visualiza os horários disponíveis;
    \item Confirma o agendamento;
    \item Recebe a confirmação do sistema.
\end{enumerate}

Esse fluxo foi projetado para simplificar o processo de marcação de consultas e reduzir etapas burocráticas.

% ---------------- TECNOLOGIAS ----------------
\section{Tecnologias Utilizadas}

O desenvolvimento do sistema SUS Mais foi realizado utilizando tecnologias modernas de desenvolvimento web, com foco em performance, escalabilidade e experiência do usuário.

As tecnologias utilizadas foram:

\begin{itemize}
    \item React 19, utilizado para construção da interface do usuário de forma componentizada e reativa;
    \item Vite, utilizado como ferramenta de build e ambiente de desenvolvimento rápido;
    \item Firebase Authentication, responsável pelo sistema de autenticação de usuários (login e cadastro);
    \item React Router, utilizado para gerenciamento de rotas e navegação entre páginas da aplicação;
    \item Bootstrap, utilizado para estilização e responsividade da interface;
    \item Axios, utilizado para realização de requisições HTTP para comunicação com APIs;
    \item React Icons, utilizado para inclusão de ícones na interface de forma prática e padronizada.
\end{itemize}

% ---------------- BENEFÍCIOS ----------------
\section{Benefícios}

A implementação do sistema SUS Mais traz diversos benefícios:

\begin{itemize}
    \item Redução de filas presenciais;
    \item Maior organização no fluxo de atendimento;
    \item Facilidade de acesso ao sistema de saúde;
    \item Otimização do tempo de usuários e profissionais;
    \item Modernização do processo de agendamento.
\end{itemize}

% ---------------- TRABALHOS FUTUROS ----------------
\section{Trabalhos Futuros}

Como melhorias futuras, o sistema pode ser expandido para:

\begin{itemize}
    \item Aplicativo mobile;
    \item Notificações automáticas por e-mail ou SMS;
    \item Integração com prontuário eletrônico;
    \item Sistema de triagem inteligente;
    \item Painel administrativo para unidades de saúde.
\end{itemize}

% ---------------- CONCLUSÃO ----------------
\section{Conclusão}

O sistema SUS Mais propõe uma solução tecnológica para modernizar o processo de agendamento de consultas no sistema público de saúde.

A proposta visa reduzir burocracias, melhorar a organização e facilitar o acesso da população aos serviços de saúde, demonstrando aplicabilidade prática e potencial de evolução para um sistema real.

\end{document}